\section{Resultados Esperados}
\label{sec:resultados_esperados}

\subsection{Entregables Técnicos}

\begin{itemize}
    \item \textbf{Sistema funcional de detección:} software completo capaz de procesar
          llamadas de call center y detectar automáticamente picos emocionales, con
          código documentado y versionado en el repositorio Git.
    \item \textbf{Flujo de trabajo reproducible:} código y cuadernos interactivos de
          Python que permitan replicar todos los experimentos, desde el
          preprocesamiento de audio hasta la detección final de picos.
    \item \textbf{Módulos independientes:} extracción de características acústicas,
          cálculo de la baseline emocional, detector de anomalías, clasificador de
          valencia y sistema de visualización.
\end{itemize}

\subsection{Métricas de Desempeño Esperadas}

\textbf{Precisión en la detección $\geq 75$\,\%:} este valor es alcanzable porque se
centra en picos emocionales extremos, con señales acústicas muy marcadas. Proyectos
similares en detección de estrés y escalamiento emocional reportan precisiones del
70--80\,\% cuando se limitan a eventos intensos \cite{abdelhamid2022, wagner2023}.

\textbf{Recall en la detección $\geq 70$\,\%:} el uso de múltiples características
acústicas permite detectar picos incluso cuando una característica individual no es
pronunciada. El baseline personalizado por hablante mejora la sensibilidad en
comparación con umbrales fijos \cite{jha2022, gat2022}.

\textbf{F1-Score $\geq 72$\,\%:} como promedio armónico entre precisión y recall,
representa un equilibrio adecuado entre evitar falsos positivos (que desperdiciarían
el tiempo de los supervisores) y capturar la mayoría de los momentos críticos.

\textbf{Clasificación de valencia $\geq 65$\,\%:} la clasificación binaria
positivo/negativo resulta más desafiante que la detección de picos, ya que las
emociones de alta activación comparten muchas características acústicas.
